% !TeX root = ../../Book2.tex
% 纲要标题：### 14.5 Burnside 矩阵代数定理
% 纲要位置：计划纲要.md，第 1235 行。

\section{Burnside 矩阵代数定理}
\label{b2:sec:1405}

代数闭域上的有限维不可约作用具有更强的结论：自同态除环只剩底域标量，因而稠密性直接给出全部矩阵。这个结论通常称为 Burnside 矩阵代数定理；它对群作用也有用，但陈述的对象是线性算子组成的代数。

% 纲要要点（供目录核对）：
%
% * 代数闭域上有限维不可约矩阵代数
% * 与稠密性定理的联系
% * 区分轨道计数引理及有限群 p^a q^b 可解性定理
%
% ---
%

\subsection{代数闭域上的不可约矩阵代数}
\label{b2:subsec:1405:01}

\begin{theorem}[Burnside 矩阵代数]\label{b2:thm:burnside-matrix-algebra}
设 \(k\) 为代数闭域，\(V\ne0\) 有限维，\(A\subseteq\operatorname{End}_k(V)\) 为含恒等算子的 \(k\) 子代数。若 \(V\) 没有非零真 \(A\) 稳定子空间，则 \(A=\operatorname{End}_k(V)\)。
\end{theorem}
\begin{proof}
假设使 \(V\) 成为简单左 \(A\) 模。由\cref{b2:cor:schur-algebraically-closed}，\(\operatorname{End}_A(V)=k\operatorname{id}_V\)：有限维与代数闭性使每个自同态有一个特征值，Schur 引理再迫使该自同态等于相应标量。因此本节的右标量除环 \(D=\operatorname{End}_A(V)^{\mathrm{op}}\) 就是 \(k\)。

\cref{b2:cor:density-finite-dimensional}于是说明作用映射 \(A\to\operatorname{End}_k(V)\) 满射。而这个映射原本就是子代数的包含，故 \(A=\operatorname{End}_k(V)\)。
\end{proof}

这一\term[Burnside-juzhen-daishu-dingli]{Burnside 矩阵代数定理}[Burnside's matrix algebra theorem]没有假设 \(k\) 的特征为零。有限维保证特征多项式有定义且次数为正，代数闭性保证它有根，再由 Schur 引理迫使上述自同态为标量。

\ProseParagraphBreak
例如，对域 \(k\) 上的一族矩阵，若其生成的含幺代数作用不可约，且 \(k\) 代数闭，则这些矩阵的有限乘积与线性组合生成全部矩阵空间。结论是代数生成，不是说原先有限个矩阵本身已经线性生成了全部矩阵。

\subsection{与稠密性定理的联系}
\label{b2:subsec:1405:02}

一般稠密性使用右除环 \(D=\operatorname{End}_A(V)^{\mathrm{op}}\)。上一定理先把它识别为 \(k\)，然后在有限基上同时指定像。代数闭假设若删去，结果可以改变：复数左乘组成 \(M_2(\mathbb R)\) 中的子代数
\[
\left\{
\vphantom{\begin{matrix}a\\b\\c\end{matrix}}
\,\begin{pmatrix}a&-b\\b&a\end{pmatrix}
\ \middle|\ a,b\in\mathbb R\,
\right\}.
\]
它在实平面上不可约，因为任何非零真稳定子空间只能是一条实直线，却不可能对乘 \(i\) 的九十度旋转稳定。但这个代数只有二维，严格小于四维的 \(M_2(\mathbb R)\)。稠密性对此给出的目标是复线性自同态，恰好就是原代数。

\ProseParagraphBreak
有限维也不能直接删除。设 \(V\) 有可数基 \(e_1,e_2,\ldots\)，\(F\) 为全部有限秩算子组成的双边理想，取 \(A=k\operatorname{id}+F\)。它含有将任意非零向量送到任意目标向量的秩一算子，所以作用不可约。可是右移算子 \(T(e_i)=e_{i+1}\) 不属于 \(A\)：对每个 \(\lambda\in k\)，\(T-\lambda\operatorname{id}\) 都单射，因为有限支集向量的最后一个非零坐标在右移后产生一个无法被抵消的末项。因此其像无限维，不可能是有限秩算子。

\ProseParagraphBreak
这个例子仍能满足有限组插值：把有限组向量所生成的子空间扩充为直和分量，在该有限维部分按要求定义算子，在一个补空间上取零，就得到所需有限秩算子。因此“在任意指定的有限组上一致”与“作为整个空间上的算子相等”确实有不同的含义。

\ProseParagraphBreak

文献中 Burnside 的名字还用于有限群作用的轨道计数引理，以及“阶为 \(p^a q^b\) 的有限群可解”的定理。本节的矩阵代数定理与这两个陈述分别讨论不同对象，不能只凭同名互相引用。特别地，矩阵代数定理没有关于群阶只含两个素因子的条件，也不直接给出群的可解性。

\ProseParagraphBreak
本定理的有限维代数闭域版本，可对照 Etingof 等人的\cite[Section 3.2, Corollary 3.2.1, Theorem 3.2.2(i), p. 43]{EtingofEtAl2011RepresentationTheory}；该书将它列为稠密性定理的一部分。后续研究有限群表示时，群代数的像正是一类矩阵代数，因此这个结果会进入表示论；本章给出的证明只依赖 Schur 引理、自同态的特征值和\cref{b2:thm:jacobson-density}。
